\subsection{Introduction}

Trong khoa học, mục tiêu cốt lõi của nghiên cứu là xây dựng và kiểm chứng các giả thuyết về thế giới thực. Theo triết gia khoa học \textbf{Karl Popper}, một lý thuyết hoặc giả thuyết khoa học phải là một phát biểu có thể được kiểm chứng bằng dữ liệu thực nghiệm. Nếu một giả thuyết không thể bị bác bỏ hoặc không thể được kiểm tra một cách hợp lý thông qua quan sát và thí nghiệm, thì nó không được xem là một giả thuyết khoa học. Quan điểm này đã đưa \textbf{dữ liệu}, \textbf{xác suất} và \textbf{thống kê} trở thành những thành phần trung tâm của phương pháp khoa học hiện đại.

Sau khi xây dựng một giả thuyết khoa học, người nghiên cứu thường chuyển nó thành các \textbf{giả thuyết thống kê} có thể kiểm định được. Khi đó, câu hỏi đặt ra là: làm thế nào để kiểm định một giả thuyết thống kê? Trong thống kê cổ điển, quá trình này có thể được xem như một dạng \textit{chứng minh phản chứng ngẫu nhiên} (\textit{stochastic proof by contradiction}). Cụ thể, ta giả định rằng \textbf{giả thuyết không} ($H_0$) là đúng, sau đó thiết kế thí nghiệm và thu thập dữ liệu. Nếu dữ liệu quan sát được quá bất thường dưới giả định $H_0$, ta bác bỏ giả thuyết không và kết luận rằng đã tìm thấy bằng chứng cho một khám phá mới.

Trong nhiều thập kỷ, bằng chứng chống lại giả thuyết không chủ yếu được đo lường bằng \textbf{\textit{p-value}}. Tuy nhiên, cách tiếp cận này thường giả định \textbf{kích thước mẫu cố định} trước khi thu thập dữ liệu. Trong khi đó, các ứng dụng hiện đại như thí nghiệm trực tuyến, thử nghiệm lâm sàng thích nghi, hệ thống học máy và phân tích dữ liệu thời gian thực thường tạo ra dữ liệu theo cách \textbf{tuần tự (sequential)}. Trong những tình huống này, việc \textbf{dừng thí nghiệm tùy ý} (\textit{optional stopping}) hoặc kiểm định lặp đi lặp lại có thể làm mất tính hợp lệ của các phương pháp suy luận truyền thống.

Để giải quyết những hạn chế trên, khuôn khổ \textbf{\textit{Sequential Anytime-Valid Inference (SAVI)}} đã được phát triển nhằm xây dựng các phương pháp suy luận vẫn duy trì được tính đúng đắn tại mọi thời điểm dừng quan sát. SAVI cho phép dữ liệu được theo dõi liên tục mà không làm gia tăng \textbf{xác suất sai lầm loại I} (\textit{Type-I error}), đồng thời cung cấp các công cụ đo lường bằng chứng có khả năng thích ứng với môi trường dữ liệu tuần tự.

Một ý tưởng trung tâm của SAVI là diễn giải việc kiểm định giả thuyết dưới góc nhìn \textbf{cá cược} (\textit{betting}) hoặc \textbf{đánh cược} (\textit{gambling}) chống lại giả thuyết không. Thay vì chỉ tính toán xác suất quan sát dữ liệu dưới giả thuyết $H_0$, ta tưởng tượng một người chơi liên tục đặt cược dựa trên dữ liệu thu được. Nếu người chơi có thể làm gia tăng tài sản của mình một cách đáng kể khi đặt cược chống lại $H_0$, thì sự tăng trưởng đó được xem là bằng chứng chống lại giả thuyết không. Theo nghĩa này, việc bác bỏ một giả thuyết tương đương với việc ``kiếm tiền'' bằng cách cá cược chống lại giả thuyết đó.

Về mặt toán học, cách tiếp cận này được xây dựng dựa trên các \textbf{\textit{test martingale}}, tức các martingale không âm bắt đầu từ giá trị 1. Từ đó xuất hiện khái niệm \textbf{\textit{e-value}} và \textbf{\textit{e-process}}, đóng vai trò như những thước đo bằng chứng thay thế cho p-value trong thống kê cổ điển. Các công cụ này tạo thành nền tảng của SAVI và cho phép thực hiện kiểm định giả thuyết hợp lệ dưới các chiến lược thu thập dữ liệu linh hoạt.

Một điểm đáng chú ý là cách nhìn nhận dựa trên cá cược tạo ra cầu nối tự nhiên giữa nhiều lĩnh vực nghiên cứu khác nhau. Từ \textbf{Lý thuyết thông tin} (\textit{Information Theory}), nơi xuất hiện các khái niệm như \textbf{entropy}, \textbf{độ phân kỳ Kullback--Leibler (KL divergence)} và \textbf{tiêu chuẩn Kelly}; đến \textbf{Xác suất và Tài chính toán học} với \textbf{martingale} và \textbf{quá trình vốn}; hay \textbf{Học trực tuyến và Tối ưu hóa} với các bài toán dự đoán tuần tự và tối thiểu hóa regret. Do đó, \textbf{\textit{Game-theoretic Statistics}} không chỉ cung cấp một nền tảng mới cho suy luận thống kê mà còn thống nhất nhiều ý tưởng quan trọng trong học máy, tài chính và lý thuyết thông tin.

Trong báo cáo này, chúng tôi sẽ lần lượt giới thiệu khái niệm \textbf{SAVI}, trình bày cách \textbf{kiểm định giả thuyết bằng cá cược}, phân tích \textbf{tiêu chuẩn Kelly} và \textbf{tính tối ưu logarit}, sau đó giới thiệu các khái niệm cốt lõi như \textbf{e-value}, \textbf{e-process} và các \textbf{chiến lược cá cược tối ưu} trong khuôn khổ \textit{Game-theoretic Statistics}. 